% Author-supplied version approved for inclusion on 1 November 2026.
\subsection*{Tools used}

\begingroup
\small
\setstretch{1.05}
\setlength{\tabcolsep}{4pt}
\renewcommand{\arraystretch}{1.2}
\begin{longtable}{@{}>{\raggedright\arraybackslash}p{.13\textwidth}>{\raggedright\arraybackslash}p{.55\textwidth}>{\raggedright\arraybackslash}p{\dimexpr.32\textwidth-4\tabcolsep\relax}@{}}
\toprule
Tool & Provider / model & Period of use\\
\midrule
\endfirsthead
\toprule
Tool & Provider / model & Period of use\\
\midrule
\endhead
\midrule
\multicolumn{3}{r}{\textit{Continued on the next page}}\\
\endfoot
\bottomrule
\endlastfoot
ChatGPT & OpenAI, ChatGPT 6 Sol & June and September 2026\\
\addlinespace
Codex & OpenAI, GPT-5.6 Sol; including the \nolinkurl{academic-humanizer} skill, version 0.3.3 & September--November 2026\\
\addlinespace
Claude & Anthropic, Opus 5.5, medium setting & October 2026\\
\end{longtable}
\endgroup

\clearpage
\subsection*{Influence on process and product}

\begingroup
\small
\setstretch{1.05}
\setlength{\tabcolsep}{4pt}
\renewcommand{\arraystretch}{1.2}
\begin{longtable}{@{}>{\raggedright\arraybackslash}p{.14\textwidth}>{\raggedright\arraybackslash}p{.10\textwidth}>{\raggedright\arraybackslash}p{.47\textwidth}>{\raggedright\arraybackslash}p{\dimexpr.29\textwidth-6\tabcolsep\relax}@{}}
\toprule
Phase & Tool & Use & Affected parts\\
\midrule
\endfirsthead
\toprule
Phase & Tool & Use & Affected parts\\
\midrule
\endhead
\midrule
\multicolumn{4}{r}{\textit{Continued on the next page}}\\
\endfoot
\bottomrule
\endlastfoot
Understanding and planning & ChatGPT, Codex & Translation, algorithm explanations, examples of sibling groups and finger trees, literature searches, architecture discussions and discussion of review findings & Research process; background and implementation explanations\\
\addlinespace
Implementation & Codex & Assistance with implementation changes, tests and validation tooling. The choice of a red--black finger tree as the efficient backend and the implementation of the sorting algorithm are my own. & Implementation and test files; qualification and archive-validation scripts\\
\addlinespace
Proof & Codex & English and LaTeX formulation of my extensions of the arguments of Fung et~al., source checking and subsequent revisions & Appendix~\ref{app:backend-argument}\\
\addlinespace
Experiments and results & Codex & Calibration and measurement tooling, inspection of saved records, reconstruction of result tables and revision of their interpretation & Experimental pipeline; Chapters~\ref{chap:methodology} and~\ref{chap:results} and related summaries\\
\addlinespace
Writing & Codex & English and LaTeX drafting and revision; use of Academic Humanizer to reduce repetitive qualifications and improve clarity & Abstract, main chapters and appendices\\
\addlinespace
Review & Claude & Technical and manuscript review, including identification of an omitted proof premise and an incomplete diagnostic table & Appendix~\ref{app:backend-argument} and the archived diagnostic indexed in Section~\ref{sec:appendix-external-review}\\
\end{longtable}
\endgroup

All measurement values reported in Chapter~\ref{chap:results} come from the
archived program runs indexed in Section~\ref{sec:appendix-current-formal}.
Codex assisted with the tooling used to check and present those results.

The following examples contain excerpts of recorded requests. English prompts
are translations of the Chinese originals; descriptions of outputs are summaries
rather than verbatim transcripts.

\subsection*{Example 1: reducing repetitive qualifications}
\textit{Codex, 31 October 2026}

\textbf{Prompt excerpt:}
\begin{quote}
\enquote{Do you think the repeated emphasis on the limits of the evidence needs changing?}
\end{quote}

\textbf{Output, summarized:}

The subsequent revision used the Academic Humanizer skill to reduce repetitive
qualifications and improve wording and transitions while retaining the relevant
assumptions and scope.

\textbf{Use:}

The revisions were incorporated into the manuscript. Checks for this revision
compared equations, mathematical environments, table contents and citations with
the preceding version and confirmed that these remained unchanged.

\subsection*{Example 2: correcting a proof premise}
\textit{Claude review, revision carried out with Codex, 20 October 2026}

\textbf{Request excerpts forwarded to Codex:}
\begin{quote}
\enquote{Add the missing premises to A.5/A.6.}

\enquote{Ensure that the bracketing condition appears explicitly in the text, and
identify its source as Fung, p.~88(b).}
\end{quote}

\textbf{Output, summarized:}

Codex checked the relevant passage in the source paper, added the bracketing
premise to Lemma~\ref{lem:proof-reentry} and explained its role in the available
group slot in Lemma~\ref{lem:proof-locality}.

\textbf{Use:}

I decided to address this finding and specified the required change and its
source in the requests above. The correction was incorporated into
Appendix~\ref{app:backend-argument}. It changed the proof exposition; the sorting
implementation was unchanged.

\subsection*{Example 3: completing the diagnostic table}
\textit{Claude review, revision carried out with Codex, 20 October 2026}

\textbf{Request excerpt forwarded to Codex:}
\begin{quote}
\enquote{Show all fifteen rows in Table 6.5, or present the fifteen combinations
of input size and nesting level.}
\end{quote}

\textbf{Output, summarized:}

Codex reconstructed the complete table from the saved measurements and revised
the accompanying interpretation. The revised presentation reports four cases
favoring the hook-free tree and eleven favoring the list backend.

\textbf{Use:}

I decided that the table should report all fifteen cases and specified the
required format. The complete table and revised interpretation were incorporated
into the results, abstract and conclusion. No new timing measurements were taken
for this revision.
